\input{../tex-inputs/latex-leseansicht-vorspann}

\section[Gustav Schwarzkopf an Arthur Schnitzler, 5. 8. 1913]{L04260 Gustav Schwarzkopf an Arthur Schnitzler, 5. 8. 1913}
\nopagebreak\mylabel{L04260v}
\rehead{ }\normalsize\beginnumbering\briefempfaengerindex{Schnitzler, Arthur@\textsc{Schnitzler, Arthur}!zzzSchwarzkopf, Gustav@\emph{von Gustav Schwarzkopf}!1913-08-051@{5. 8. 1913}|(be}
\toendnotes[C]{\smallbreak\pagebreak[2]}
\correspDesc{Versand  durch Gustav Schwarzkopf am 5. 8. 1913 in Wien
\newline{}Erhalt  durch Arthur Schnitzler im Zeitraum [6. 8. 1913 –
                  7. 8. 1913?] in Brijuni}\toendnotes[C]{\smallbreak}
\Standort{CUL, Schnitzler, B 96a.}
\physDesc{Briefkarte, 1049 Zeichen
\newline{}Handschrift: schwarze Tinte, deutsche Kurrent
\newline{}Schnitzler: mit Bleistift Vermerk: »\textsc{Schwarzk}{[}opf{]}« }\toendnotes[C]{\smallbreak}
\pstart
           {\pb}\textsc{Wien\oindex{Wien@\textbf{Wien}, \emph{Verwaltungsgebiet}|pw}}{ }5. VIII 13\pend
           \vspace{0.5em}
\pstart
           Lieber Arthur, Sie werden{ }ſich vielleicht gewundert haben — oder
               vielleicht auch nicht – daß ich gar nichts von mir hören laſſe, nicht einmal meine
               Ankunft melde. Aber \uline{dies}mal war ich wirklich
               verhindert, weil man mich 24 Stunden nach meiner Ankunft in’s Sanatorium »\textsc{Hera}\oindex{Wien@\textbf{Wien}!IX., Alsergrund@\textbf{IX., Alsergrund}!Sanatorium Hera@\textbf{Sanatorium Hera}, \emph{Sanatorium}|pw}« gebracht hat – aber wirklich \uline{gebracht} mit
               Tragbahre! Sie ſehen, wie{ }ſtolz mich das macht – als »blinddarmverdächtig«. Uns 4 ½
               Tagen »Beobachtung« hat man mich wieder geſund entlaſſen, ohne zu ſchneiden. Das
               Intermezzo war ſchmerzhaft, dafür aber nicht billig. {\pb}Das vornehmſte Zeichen meiner Erholung –
               die blühende Naſe – habe ich dabei ſo ziemlich eingebüßt. Mein Bruder Max\pwindex{Schwarzkopf, Max 12.\,6.\,1857 Wien – 14.\,4.\,1928 ebd.@\textsc{Schwarzkopf, Max} (12.\,6.\,1857 Wien – 14.\,4.\,1928 ebd.), \emph{Rechtsanwalt}|pw} wollte gerade abreiſen, als das Ereigens
               eintrat und iſt heute noch nicht fort; \uline{jetzt} wartet
               er auf den erſten leidlich hübſchen Tag.\pend
           
\pstart
           Hoffentlich haben Sie in Brioni\oindex{Brijuni@\textbf{Brijuni}|pw} beſſeres
               Wetter. Wie gut, daß ich nicht geblieben bin,{ }ſonſt hätte ich ihnen dieſe Scherereien
               gemacht.\pend
           
\pstart
           Ihre Frau\pwindex{Schnitzler, Olga 17.\,1.\,1882 Wien – 13.\,1.\,1970 Lugano@\textsc{Schnitzler, Olga} (17.\,1.\,1882 Wien – 13.\,1.\,1970 Lugano), \emph{Schauspielerin, Sängerin}|pwv}, Sie und die
                  Kinder\pwindex{Schnitzler, Heinrich 9.\,8.\,1902 Hinterbrühl – 12.\,7.\,1982 Wien@\textsc{Schnitzler, Heinrich} (9.\,8.\,1902 Hinterbrühl – 12.\,7.\,1982 Wien), \emph{Regisseur, Schauspieler}|pwv}\pwindex{Cappellini, Lili 13.\,9.\,1909 Wien – 26.\,7.\,1928 Venedig@\textsc{Cappellini, Lili} (13.\,9.\,1909 Wien – 26.\,7.\,1928 Venedig)|pwv}
               herzlichſt grüßend{\\[\baselineskip]}Ihr{\\[\baselineskip]}\spacefill\mbox{Gustav.}\pend
           \leftskip=0em{}
\pstart
           \noindent{}Ich erteile ihnen Vollmacht mich ſä{\geminationm}tlichen Beka{\geminationn}ten beſtens zu empfehlen.\pend
           \selectlanguage{ngerman}\endnumbering\briefempfaengerindex{Schnitzler, Arthur@\textsc{Schnitzler, Arthur}!zzzSchwarzkopf, Gustav@\emph{von Gustav Schwarzkopf}!1913-08-051@{5. 8. 1913}|)be}\mylabel{L04260h}
\begin{anhang}
\end{anhang}\newcommand{\dateiname}{L04260}\newcommand{\titel}{Gustav Schwarzkopf an Arthur Schnitzler, 5. 8. 1913}\newcommand{\editorInnen}{Herausgegeben von Jahnke, SelmaMüller, Martin Anton}\input{../tex-inputs/latex-leseansicht-abspann}
